%%%%%%%%%%%%%%%%%%%%%%%%%%%%%%%%%%%%%%%%%%%%%%%%%%%%%%%
% % Lines starting with % are comments, which are ignored.
% % This is a handy way of indicating the date and version of
% % your document, to wit:
% %
% % LaTeX sample file
% % Modified March, 2002
% %
%%%%%%%%%%%%%%%%%%%%%%%%%%%%%%%%%%%%%%%%%%%%%%%%%%%%%%%
\DocumentMetadata{
  % PDF standards: ONE key whose value is a comma-separated list --
  % UA-2 (accessibility) plus A-4f (archival; the archival companion of
  % UA-2). Do NOT use separate pdfstandard=... keys: with ordinary
  % key-value parsing only the last one takes effect.
  pdfstandard={UA-2,A-4f},
  lang=en-US,
  tagging=on,
  tagging-setup={
    math/setup=mathml-SE
  }
}

\documentclass{ltx-talk}
% lecture-macros.tex -- shared package setup for lecture-template.tex
% (ltx-talk slides / handouts).
%
% NOTE ON PATHS: every template in this repository is stored one level up
% from where it is compiled (see README.md). All relative \input paths
% below are resolved against the folder the document is compiled in,
% which sits one directory above this file.

% --- Tables, bookmarks, strikethrough ------------------------------------
% booktabs: professional table rules (\toprule, \midrule, \bottomrule)
% bookmark: PDF outline entries (PDF bookmarks)
% ulem[normalem]: \sout strikethrough that survives math mode
\usepackage{booktabs}
\usepackage{bookmark}
\usepackage[normalem]{ulem}
\usepackage{multicol}

% --- Math ------------------------------------------------------------
% amsmath: equation environments and numbering
% mathrsfs: \mathfrak  |  stmaryrd: extra symbols
\usepackage{amsmath}
\usepackage{mathrsfs}
\usepackage{stmaryrd}

% --- Layout -------------------------------------------------------------
% float: [H] float placement
% fontspec: Unicode/OpenType fonts (compile with XeLaTeX or LuaLaTeX)
\usepackage{float}
\usepackage{fontspec}

% --- Colors ---------------------------------------------------------------
% dvipsnames: extended named color palette (the Accessible* colors below
% are built on it);
% table: \rowcolors / \cellcolor for shading tables.
% NOTE: xcolor may only be loaded ONCE per document -- a second
% \usepackage{xcolor} with different options is a fatal "Option clash"
% error, so ALL of its options are listed here.
\usepackage[dvipsnames,table]{xcolor}

% --- Multimedia -----------------------------------------------------------
% animate: embeds audio/video (and legacy Flash) media in a frame.
\usepackage{animate}

% --- Bibliography ----------------------------------------------------------
% biblatex (biber backend): citations for lecture notes.
\usepackage[
  backend=biber,
  style=alphabetic,
  url=true,
  doi=true,
  isbn=true,
  backref=true
]{biblatex}

% --- Colored boxes ----------------------------------------------------------
% talkthemetcolorbox: the colored-box theme for ltx-talk slides
% (shadow/rounded/shaded blocks, showtitle, titlepage styles).
\usepackage[
  shadow,
  rounded,
  shaded,
  blocks,
  showtitle,
  titlepage
]{talkthemetcolorbox}

% --- "theindex" helper environment ------------------------------------------
% Provides a lightweight theindex environment (an itemize) and the
% \@idxitem/\subitem/\subsubitem/\indexspace macros that LaTeX's standard
% theindex machinery expects, so hand-written index pages indent correctly.
\makeatletter
\newenvironment{theindex}
  {\begin{itemize}}
  {\end{itemize}}
\providecommand\@idxitem{\par\hangindent 40\p@}
\providecommand\subitem{\@idxitem \hspace*{20\p@}}
\providecommand\subsubitem{\@idxitem \hspace*{30\p@}}
\providecommand\indexspace{\par \vskip 10\p@ \@plus5\p@ \@minus3\p@\relax}
\makeatother

% --- Index -----------------------------------------------------------------
% imakeidx: builds the printed index.
% The redefinition below appends "|hyperpage" to every index entry so that,
% in the PDF, each index entry is a link to its page.
\usepackage{imakeidx}
\makeatletter
\let\imki@orig@wrindexentrysplit\imki@wrindexentrysplit
\renewcommand\imki@wrindexentrysplit[3]{%
  \imki@orig@wrindexentrysplit{#1}{#2|hyperpage}{#3}%
}

\makeatletter
% --- Section titles: capture + one bookmark per logical section -------------
% \section is redefined so that:
%   1. the section title is captured into the plain-TeX register
%      \mytalksectitletoks (reprinted verbatim on the auto divider frame,
%      see the cmd/section/after hook below -- plain-TeX storage is
%      catcode-safe, unlike expanding the title in a hook), and
%   2. exactly ONE depth-1 PDF bookmark is created per \section call
%      (hyperref's own bookmark is suppressed via \Hy@writebookmark).
\DeclareCommandCopy{\OriginalSectionForBookmarks}{\section}
\newtoks\mytalksectitletoks

\RenewDocumentCommand{\section}{s o m}{%
  \begingroup
    \def\Hy@writebookmark##1##2##3##4##5{}%
    \global\mytalksectitletoks={#3}%
    \IfBooleanTF{#1}
      {\phantomsection\IfNoValueTF{#2}{\OriginalSectionForBookmarks*{#3}}{\OriginalSectionForBookmarks*[#2]{#3}}}%
      {\IfNoValueTF{#2}{\OriginalSectionForBookmarks{#3}}{\OriginalSectionForBookmarks[#2]{#3}}}%
  \endgroup
  \IfBooleanTF{#1}
    {\IfNoValueTF{#2}{\pdfbookmark[1]{#3}{section-star.\thepage.\number\inputlineno}}{\pdfbookmark[1]{#2}{section-star.\thepage.\number\inputlineno}}}%
    {\IfNoValueTF{#2}{\pdfbookmark[1]{Section \thesection: #3}{section.\thesection}}{\pdfbookmark[1]{Section \thesection: #2}{section.\thesection}}}%
}
\makeatother

% --- One PDF bookmark per logical frame (NOT per overlay) ---------------------
% ltx-talk writes one PDF page per overlay, so a naive frametitle hook would
% register a new bookmark for every overlay of the same frame. The
% \g_slideutils_frame_bookmark_int register remembers the last frame that
% got a bookmark, so subsequent overlays of the same frame are skipped.
% This is also what keeps BeamerPresenter's cumulative annotations working:
% one logical frame must stay one logical unit in the PDF outline.
\ExplSyntaxOn
\int_new:N \g_slideutils_frame_bookmark_int

\AddToHook{cmd/frametitle/after}{%
  \tl_if_empty:NF \g__talk_frame_title_tl
    {
      \ifnum\value{frame}=\g_slideutils_frame_bookmark_int\relax
      \else
        \int_gset:Nn \g_slideutils_frame_bookmark_int { \value{frame} }%
        \pdfbookmark[2]{\tl_use:N \g__talk_frame_title_tl}
          {frame.\int_use:N \g_slideutils_frame_bookmark_int}
      \fi
    }
}
\ExplSyntaxOff
\makeatother

% --- Printed index + theme for the colored boxes -------------------------------
% \makeindex[...]: two-column index titled "Index", listed in the TOC.
% The two EditInstance calls style the standard header/footer instances
% of talkthemetcolorbox (cyan background; footer order title/author/
% date/framenumber).
\makeindex[columns=2, title=Index, intoc]\EditInstance{header}{std}{
  background-color = cyan,
  color            = black,
  height           = 0.75cm
}
\EditInstance{footer}{std}{
  background-color = cyan,
  element-order = {title, author, date, framenumber},
}

% --- Auto section divider frame ---------------------------------------------------
% After each \section, insert a centered, large, structure-colored frame
% that prints the section title (from \mytalksectitletoks).
\AddToHook{cmd/section/after}{%
  \begin{frame}
    \centering
    \color{structure}
    \LARGE
    \the\mytalksectitletoks
  \end{frame}%
}

% --- Tagging sockets (tagged PDF structure) ---------------------------------------
% ltx-talk exposes tagging "sockets" where the PDF structure can attach
% section titles; register one and leave it unassigned (none) so the
% class's own tagging keeps control.
\ExplSyntaxOn
\AtBeginDocument{
  \NewTaggingSocketPlug { talk / sec / title } { none } { }
  \AssignTaggingSocketPlug { talk / sec / title } { none }
}
\ExplSyntaxOff

% --- Accessibility roles ----------------------------------------------------------------
% Map \frametitle to the H1 heading role in the tagged PDF, so screen
% readers treat slide titles as headings.
\tagpdfsetup{role / new-tag = frametitle / H1}

\makeatletter
% --- Page label = frame number -------------------------------------------------------
% Redefine what feeds the PDF page label to track the frame counter,
% not the physical page counter. Replace \theframenumber below with
% whatever internal macro/property ltx-talk actually uses to produce
% the "framenumber" footer element — confirm via ltx-talk-frame.dtx.
\renewcommand\thepage{\arabic{frame}}
\makeatother

\addbibresource{refs.bib}

% --- "Problem Time!" slides ------------------------------------------------------------------
% \ProblemSlide{problem}{solution} emits two frames: the problem, then
% its solution after a \pause, each numbered by the problem counter.
\newcounter{problemcounter}
\setcounter{problemcounter}{0}
\newcommand{\ProblemSlide}[2]{
  \stepcounter{problemcounter}
  \begin{frame}
    \frametitle{Problem Time!}
    Problem \theproblemcounter: #1
  \end{frame}

  \begin{frame}
    \frametitle{Problem Time!}
    Solution to Problem \theproblemcounter: \pause

    #2
  \end{frame}
}
% \highlight[<colour>]{<stuff>}
% from https://tex.stackexchange.com/questions/33401/a-version-of-colorbox-that-works-inside-math-environments
\newcommand{\highlight}[2][yellow]{\mathchoice%
  {\colorbox{#1}{$\displaystyle#2$}}%
  {\colorbox{#1}{$\textstyle#2$}}%
  {\colorbox{#1}{$\scriptstyle#2$}}%
  {\colorbox{#1}{$\scriptscriptstyle#2$}}}%

% --- Fonts ------------------------------------------------------------------------------
% Lato (text) + LeteSansMath (math): open, accessible typefaces.
\setmainfont{Lato}[
  Ligatures=TeX,
  Numbers={Lining,Monospaced}
]
\setsansfont{Lato}[
  Ligatures=TeX,
  Numbers={Lining,Monospaced}
]
\setmathfont{LeteSansMath.otf}\usepackage{emoji}

% --- Accessible colors ---------------------------------------------------------------------
% WCAG-contrast-safe variants of red/green for text on white:
% red!80!black and green!40!black both clear the 4.5:1 requirement.
% AltYellow: soft background for highlight boxes.
\colorlet{AccessibleRed}{red!80!black}
\colorlet{AccessibleGreen}{green!40!black}
\definecolor{AltYellow}{RGB}{255,245,170}

% --- Links ---------------------------------------------------------------------------------
% hyperref + link colors: every link uses AccessibleGreen (7.2:1 on
% white), so internal, citation, and URL links all pass WCAG 1.4.3.
\usepackage{hyperref}
\hypersetup{
  colorlinks=true,
  linkcolor=AccessibleGreen,
  urlcolor=AccessibleGreen,
  citecolor=AccessibleGreen,
}

% --- Shared course macros ---------------------------------------------------------------------
% Root macros.tex: math shorthands used across all course documents
% (\RR, \NN, \ZZ, \QQ, \grad, \curl, \tr, \der, \pder, ...).
\newcommand{\RR}{\mathbb{R}}
\newcommand{\NN}{\mathbb{N}}
\newcommand{\ZZ}{\mathbb{Z}}
\newcommand{\QQ}{\mathbb{Q}}
\newcommand{\vv}{\vec{v}}
\newcommand{\ww}{\vec{w}}
\newcommand{\xx}{\vec{x}}
\newcommand{\yy}{\vec{y}}
\newcommand{\zz}{\vec{z}}
\newcommand{\uu}{\vec{u}}
\newcommand{\basis}[1]{\vec{e_{#1}}}
\newcommand{\id}{\mathrm{id}}
\newcommand{\CC}{\mathbb{C}}
\renewcommand{\div}{\operatorname{div}}
\newcommand{\grad}{\operatorname{grad}}
\newcommand{\curl}{\operatorname{curl}}
\newcommand{\tr}{\operatorname{tr}}
\newcommand{\der}[2]{\dfrac{d #1}{d #2}}
\newcommand{\pder}[3][]{\dfrac{\partial^{#1} #2}{\partial #3}}


% % Title and author(s)
%%%%%%%%%%%%%%%%%%%%%%%%%%%%%%%%%%%%%%%%%%%%%%%%%%%%%%%
\title{Lecture 8 - Continuity and Differentiability}
\author{Garrett Credi}
\date{10/01/2026}
%%%%%%%%%%%%%%%%%%%%%%%%%%%%%%%%%%%%%%%%%%%%%%%%%%%%%%%
% %
% % The next command allows your in import encapsulated
% % postscript files, .epsf or .eps files, which
% % contain vector graphic image data.
% %
%%%%%%%%%%%%%%%%%%%%%%%%%%%%%%%%%%%%%%%%%%%%%%%%%%%%%%%
% %   The first thanks indicates your affiliation
% %
% %  Just the name here.
% %
% % Your mailing address goes at the end.
% %
% % \thanks is also how you indicate grant support
% %
%%%%%%%%%%%%%%%%%%%%%%%%%%%%%%%%%%%%%%%%%%%%%%%%%%%%%%%

%%%%%%%%%%%%%%%%%%%%%%%%%%%%%%%%%%%%%%%%%%%%%%%%%%%%%%%
% $ The next command ensures the margins are 1 inch on
% % all sides.
% %
%%%%%%%%%%%%%%%%%%%%%%%%%%%%%%%%%%%%%%%%%%%%%%%%%%%%%%%
\begin{document}
\begin{frame}
  \maketitle
\end{frame}

\begin{frame}
  \frametitle{Announcements}

  \begin{enumerate}
    \item Exam 1 is next Tuesday (10/06/2026) in class. 
    \item Exam 1 Review Session is next Monday (10/05/2026) from 6:10-8:10 PM in Vincent Hall 1.
    \item All quizzes will be graded by Friday Night. If you want me to look at your homeworks, please email me and I'll grade them specifically.
    \item Practice Exam Solutions will be released after class today.
    \item Compiled lecture slides also released after class today.
  \end{enumerate}
\end{frame}

\begin{frame}
  \frametitle{Quiz Time}

  Quiz Time! Let's get into neat rows and spread out - only then will I pass out the quizzes.

  \textcolor{red}{DO NOT START OR FLIP OVER THE QUIZ UNTIL I OFFICIALLY START THE QUIZ.} If you have any questions about the rules for quizzes, refer to the slide in Lecture 2 or ask me directly. No cheating as per usual.
\end{frame}

\section{Continuity}

\begin{frame}
  \frametitle{Continuity Recap I}

  Recall our definitions of continuity for a function $f(x)$:

  \begin{enumerate}
    \item $f(x)$ is continuous at $x = a$ if: \pause
      \vspace{.1in}
    \item $f(x)$ is continuous on $(a,b)$ if: \pause
            \vspace{.1in}
    \item $f(x)$ is continuous \textit{everywhere} if: \index{continuity!everywhere} \pause
            \vspace{.1in}
    \item $f(x)$ is continuous if: \pause
            \vspace{.1in}
  \end{enumerate} \pause
\end{frame}
\begin{frame}
  \frametitle{Continuity Recap II}
  Also recall that: if $f(x)$ and $g(x)$ are continuous at $x = a$ and $h(x)$ is continuous at $f(a)$ then:
  \begin{multicols}{2}
    \begin{itemize}
      \item $f(x) + g(x)$ is continuous at $x = a$
      \item $f(x) - g(x)$ is continuous at $x = a$
      \item $f(x) \cdot g(x)$ is continuous at $x = a$
      \item If $g(a) \neq 0$, then $\dfrac{f(x)}{g(x)}$ is continuous at $x = a$
      \item $h(f(x))$ is continuous at $x = a$.
    \end{itemize}
  \end{multicols}
\end{frame}

\begin{frame}
  \frametitle{Continuity Recap III}
  Finally we determined a large family of continuous functions:
  \begin{multicols}{2}
    \begin{itemize}
      \item Constant functions $f_c(x) = c$
      \item Identity function $f(x) = x$
      \item Polynomials $p(x) = a_nx^n + a_{n-1}x^{n-1} + \ldots + a_1 x + a_0$
      \item Root functions $f(x) = \sqrt[n]{x}$
      \item Rational functions $\dfrac{p(x)}{q(x)}$
      \item Trigonometric functions $\sin(x), \cos(x), \tan(x), \arcsin(x)$ etc.
      \item Exponential functions $f(x) = c^x$
      \item Logarithmic functions $f(x) = \log_c(x)$
    \end{itemize}
  \end{multicols}
\end{frame}

\begin{frame}
  \frametitle{Classifying Failures of Continuity: Jump Discontinuities}
  What does discontinuity look like?

  \begin{block}{Jump Discontinuity}\index{continuity!jump discontinuity}\index{jump discontinuity}\index{discontinuity!jump discontinuity}
    A function $f(x)$ has a \textit{jump discontinuity} at $x = a$,
    \pause if the value of $f(x)$ jumps as we pass over $x = a$. \pause

    Explicitly: \pause
    \begin{enumerate}
      \item The function limits towards a value from the left: \pause
        $\lim\limits_{x \rightarrow a^{-}} f(x)$ exists \pause
      \item The function limits towards a value from the right: \pause
        $\lim\limits_{x \rightarrow a^{+}} f(x)$ exists \pause
      \item Those values do not agree: \pause
        \begin{equation}
          \lim\limits_{x \rightarrow a^{-}} f(x) \neq \lim\limits_{x \rightarrow a^{+}} f(x)
        \end{equation}
    \end{enumerate}
  \end{block}
\end{frame}

\begin{frame}
  \frametitle{Classifying Failures of Continuity: Jump Discontinuities Visually}

  Visually, jump discontinuities look like:
  \vfill
\end{frame}

\begin{frame}
  \frametitle{Classifying Failures of Continuity: Removable Discontinuities}
  Best kind of discontinuity. \pause
  \begin{block}{Removable Discontinuity}\index{continuity!removable discontinuity}\index{removable discontinuity}\index{discontinuity!removable discontinuity}\index{continuity!hole}\index{hole}\index{discontinuity!hole}
    A function $f(x)$ has a \textit{removable discontinuity} at $x = a$ if: \pause
    \begin{enumerate}
      \item $\lim\limits_{x \rightarrow a} f(x)$ exists \pause
      \item $f(a)$ either does not exist \pause (we specifically call this a \textit{hole discontinuity} or just a \textit{hole}) \pause or is not equal to $\lim\limits_{x \rightarrow a} f(x)$.
    \end{enumerate}
  \end{block} \pause

  Can redefine $f(x)$ \textit{at a single point} to make it continuous at $x = a$.
\end{frame}

\begin{frame}
  \frametitle{Removable Discontinuities Visually}
  Visually,
  \vfill
\end{frame}

\begin{frame}
  \frametitle{Classifying Failures of Continuity: Infinite Discontinuities}
  \begin{block}{Infinite Discontinuity}\index{continuity!infinite discontinuity}\index{infinite discontinuity}\index{discontinuity!infinite discontinuity}
    A function $f(x)$ has an \textit{infinite discontinuity} if, equivalently: \pause
    \begin{itemize}
      \item $f(x)$ has a vertical asymptote at $x = a$ or \pause
    \item $\lim_{x \rightarrow a} |f(x)| = + \infty$
    \end{itemize}
  \end{block}
\end{frame}

\begin{frame}
  \frametitle{Infinite Discontinuities Visually}
  Visually, 
  \vfill
\end{frame}

\begin{frame}
  \frametitle{Classifying Failures of Continuity: Essential Discontinuities}
  The \textit{WORST} kinds of discontinuities. \pause

  \begin{block}{Essential Discontinuity}
    A function $f(x)$ has an \textit{essential discontinuity} at $x = a$ if equivalently \pause
    \begin{itemize}
      \item At least one of the one-sided limits do not exist (not even diverging to $\pm infty$) or \pause
      \item As $x \rightarrow a$, $f(x)$ takes on every number in some interval $(v_{min}, v_{max})$ \textit{infinitely many times}
    \end{itemize}
  \end{block}
\end{frame}

\begin{frame}
  \frametitle{Essential Discontinuities Visually}
  Visually,
  \vfill
\end{frame}

\begin{frame}
  \frametitle{What does Continuity buy us?}

  We already saw one important use of continuous functions: \pause

  \begin{alertblock}{Direct Substitution}\index{continuity!limit direct substitution}\index{Limit Direct Substitution Theorem}
    If $f(x)$ is continuous at $x = a$ then:
    \begin{equation}
      \lim_{x \rightarrow a} f(x) = f(a)
    \end{equation} \pause
  \end{alertblock}

  Continuous functions also have another very important property: \pause

  \begin{alertblock}{Intermediate Value Theorem}\index{continuity!intermediate value theorem}\index{intermediate value theorem}
    If $f(x)$ is continuous on $(a,b)$ then \pause
    (assuming $f(a) \leq f(b)$) \pause 

    for any intermediate output $c$ so that $f(a) \leq c \leq f(b)$, there must an input $x_0$ so that $a \leq x_0 \leq b$ and $f(x_0) = c$.
  \end{alertblock} \pause

  Tagline: all intermediate output values must be hit by intermediate input values.
\end{frame}

\begin{frame}
  \frametitle{IVT Example: Root finding I}
  Question: For the function
  \begin{equation}
    f(x) = x^3 - 3x^2 + 2x + 1,
  \end{equation}
  which of the following intervals must contain a root of $f(x)$?
  \begin{multicols}{2}
    \begin{enumerate}
      \item $[-2, -1]$
      \item $[-1, 0]$
      \item $[0, 1]$
      \item $[1, 2]$
    \end{enumerate}
  \end{multicols} \pause

  \vspace{.4in}
  Approach 1: Solving. \pause IDK how to factor general cubics! \pause

  Approach 2: IVT. \pause We want to find an interval which contains a value $x_0$ so that $f(x_0) = 0$. \pause Note: if $a < 0$ is negative and $b > 0$ is positive, \pause then $0$ is always inside $[a, b]$. \pause
\end{frame}

\begin{frame}
  \frametitle{IVT Example: Root Finding II}

  Goal: find which of the intervals
    \begin{multicols}{2}
    \begin{enumerate}
      \item $[-2, -1]$
      \item $[-1, 0]$
      \item $[0, 1]$
      \item $[1, 2]$
    \end{enumerate}
  \end{multicols} \pause
  makes
  \begin{equation}
    f(x) = x^3 - 3x^2 + 2x + 1
  \end{equation}
  switch signs at the endpoints. \pause

  \begin{multicols}{2}
    \begin{enumerate}
      \item $f(-2) = -8 - 3 \cdot 4 - 2 \cdot 2 + 1 = -8 - 12 - 4 + 1 = -23$ \pause
      \item $f(-1) = -1 - 3 - 2 + 1 = -5$ \pause
      \item $f(0) = 1$ \pause ROOT! \pause
      \item $f(1) = 1 - 3 + 2 + 1 = 1$ \pause
      \item $f(2) = 8 - 12 + 4 + 1 = 1$ \pause
    \end{enumerate}
  \end{multicols}

  Only know that a root must be inside $[-1, 0]$.
\end{frame}

\section{Differentiability}

\begin{frame}
  \frametitle{Differentiability Recap}

  Again, recall that: \pause
  \begin{enumerate}
    \item $f(x)$ is differentiable at $x = a$: 
    \item $f(x)$ is differentiable on $(a, b)$:
    \item $f(x)$ is differentiable everywhere: \index{differentiability!everywhere} 
    \item $f(x)$ is differentiable
  \end{enumerate}
\end{frame}

\begin{frame}
  \frametitle{Classifying Failures of Differentiability: Discontinuities}

  We can already classify some behavior making functions differentiable. \pause

  Recall: if $f(x)$ is not \textit{continuous} at $x = a$ then it cannot be \textit{differentiable} at $x = a$. 
  \vfill
\end{frame}

\begin{frame}
  \frametitle{Classifying Failures of Differentiability: Corners}
  We can classify the rest of the behavior by examining the \textit{continuity} of the derivative function $f'(x)$: \pause

  \begin{block}{Corner Singularity}\index{differentiability!corner singularity}\index{corner singularity}\index{singularity!corner}
    For a function $f(x)$ continuous at $x = a$, we say $f(x)$ has a \textit{corner singularity} or \textit{corner} at $x = a$ \pause
    if the derivative function $f'(x)$ has a \textit{jump discontinuity} at $x = a$.
  \end{block} \pause
  Visually,
  \vfill
\end{frame}

\begin{frame}
  \frametitle{Classifying Failures of Differentiability: Cusps}
  Similarly, \pause
  \begin{block}{Cuspidal Singularity}\index{differentiability!cuspidal singularity}\index{cusp}\index{cuspidal singularity}\index{singularity!cuspidal}
    For a function $f(x)$ continuous at $x = a$, we say $f(x)$ has a \textit{cuspidal singularity} or \textit{cusp} at $x = a$ \pause
    if the derivative function $f'(x)$ has an \textit{infinite discontinuity} at $x = a$ \pause \textit{AND} does not diverge to $\pm \infty$. \pause (e.g. the one sided limits separately go to $+\infty$ and $-\infty$). \pause
  \end{block}

  Visually,
  \vfill
\end{frame}

\begin{frame}
  \frametitle{Classifying Failures of Differentiability: Vertical Tangents}
  \begin{block}{Vertical Tangent Singularity}\index{differentiability!vertical tangent singularity}\index{vertical tangent}\index{singularity!vertical tangent}
    For a function $f(x)$ continuous at $x = a$, we say $f(x)$ has a \textit{vertical tangent singularity} or a \textit{vertical tangent} at $x = a$ \pause
    if:
    \begin{equation}
      \lim_{x \rightarrow a} |f'(x)| = + \infty.
    \end{equation}
  \end{block} \pause
  Visually, 
  \vfill
\end{frame}

\begin{frame}
  \frametitle{Classifying Failures of Differentiability: Endpoint Singularity}
  Kind of a non-problem singularity: \pause
  \begin{block}{Endpoint Singularity}\index{endpoint singularity}\index{differentiability!endpoint singularity}\index{singularity!endpoint}
    For a function $f(x)$ defined on a \textit{closed} interval $[a,b]$, \pause
    we say that $f(x)$ has an \textit{endpoint singularity} at the endpoints
    $x = a$ and $x = b$ \pause just because two sided limits do not make sense at those points.
  \end{block} \pause
  Visually, 
  \vfill
\end{frame}

\begin{frame}
  \frametitle{Continuous Differentiability I}
  What about the last kind of discontinuity: \pause removable discontinuities?

  The last thing that could go wrong would be if $f'(x)$ had a hole
  discontinuity at $x = a$ \pause - not often discussed as a class of singularity. \pause

  Instead, we care when everything goes \textit{right}: \pause

  \begin{block}{Continuous Differentiability}\index{differentiability!continuous differentiability}\index{continuous differentiability}
    A function $f(x)$ is \textit{continuously differentiable} at $x = a$ if \pause $f'(x)$ is continuous at $x = a$. \pause

    Similarly, continuously differentiable on $(a,b)$, \pause, everywhere, etc. \pause

    A function $f(x)$ is \textit{continuously differentiable} or \textit{of class $C^1$} if is is continuously differentiable at every point in its domain.
  \end{block} \pause

  Intuituion: Functions whose graphs are continuous, and corner/cusp/vertical tangent free.
\end{frame}

\begin{frame}
  \frametitle{What does Continuous Differentiability buy us? I}
  Since $f(x)$ being continuously differentiable implies $f'(x)$ is continuous, we get the following by applying IVT to $f'(x)$: \pause
  \begin{alertblock}{Rolle's Theorem}\index{Rolle's Theorem}
    For $f(x)$ continuously differentiable, and $x_0$ and $x_1$ two roots of $f(x)$ so that $x_0 \leq x_1$, \pause

    there always exists an $a$ such that $x_0 \leq a \leq x_1$ and $f'(a) = 0$.
  \end{alertblock}
  \vfill
\end{frame}
\begin{frame}
  \frametitle{What does Continuous Differentiability buy us? II}
  This also gives us a \textit{VERY} important theorem: \pause
  \begin{alertblock}{Mean Value Theorem}\index{Mean Value Theorem}
    For $f(x)$ continuously differentiable, and $(x_0, f(x_0)), (x_1, f(x_1))$ so that $x_0 < x_1$ there is always an $a$ so that $x_0 \leq a \leq x_1$ and, equivalently: \pause
    \begin{itemize}
    \item The tangent line $y = f'(a)(x - a) + f(a)$ is parallel to the secant line $y = \dfrac{f(x_1) - f(x_0)}{x_1 - x_0}(x - x_0) + f(x_0)$ or \pause
    \item 
      \begin{equation}
        \dfrac{f(x_1) - f(x_0)}{x_1 - x_0} = f'(a).
      \end{equation}
    \end{itemize}
  \end{alertblock}
\end{frame}

\begin{frame}
  \frametitle{Uses of the Mean Value Theorem}

  \vfill
  Why care about MVT? \pause

  \begin{enumerate}
    \item If $f'(x)$ is the zero function $f'(x) = 0$, then $f(x)$ must be constant! \pause
    \vfill
    \item Extremely useful for doing derivations involving derivatives: e.g. determining quadratic approximations! \pause
    \vfill
  \item Lots of the really useful tools in multivariable calculus need to use MVT to even work.
    \vfill
  \end{enumerate}

  ``Fun" application: 
  \vfill
\end{frame}

\begin{frame}
  \frametitle{Worksheet Time}

  Let's get into groups and work on the worksheet!
\end{frame}

\begin{frame}
  \frametitle{End of Class!}

  That's all folks!
\end{frame}
\end{document}
